\scanpage{225}
\mspage{171}
% archive sheet number: 224

% MS heading
\uline{Nouvelle version de la théorie de} $\mathrm{HOT}_W$

% MS lines 2–10, opening paragraph
Je me suis aperçu que l'hypothèse $W\supset W_\infty$ (L 3 \unsure{grandes esp.|gr. des esp.|gr. des sph.!}) que j'avais fini par mettre sur $W$, \inserted{est artificielle,} et \deleted{que} l'hypothèse L 5 (p. 43) voire L 6, L 6 bis (p. 43-44), est inutile pour une grande partie de la théorie de $\mathrm{HOT}_W$. J'ai été ainsi amené à revoir la théorie (i.e. ce que j'en ai fait jusqu'à présent) sous les hypothèses suivantes sur $W$ (cf $\overline{\underline{\mathrm{XII}}}$) :

% MS lines 11–13, three conditions
$W(1)$ (\uline{Saturation} de $W$) : comme avant (L 1 page 1)

$W(2)$ (localisation) : c'est l'ancien L 4 bis$^*$

$W(3)$ (objet final) : Si $X$ a un objet final, $X$ est $W$-asphérique (ancien L 3)

% MS left margin against W(2)–W(3)
% slanted
\margin{i.e. \deleted{\unsure{L5|le}} le L 5 de la page 2 (l'axiome est repris (\unsure{sous une|en une} forme voisine et \unsure{équivalente}) \unsure{au} \deleted{dans} de la page 43.}

% MS lines 14–18, sufficiency
Ces \deleted{\unsure{quatre}} \deleted{\unsure{trois}} axiomes sont suffisants pour tout ce que j'avais fait dans P.S., et dans bien grande partie de ce qui a été fait ici sur HOT. Parfois cependant il faut$^*$ l'ancien L 5, (que je réénonce ici (voyez $\overline{\underline{\mathrm{XII}}}$ th. )) :

% MS lines 19–21, relative condition
$W(4\,\mathrm b\unsure{{}^*|{}'})$ : Soient $X,Y$ $W$-fibrés sur $S$, $f:X\longrightarrow Y$ un $S$-morphisme. Si $f\in W$, alors $\forall s\in\mathrm{Ob}\,S$, $f_s:X_s\longrightarrow Y_s$ est dans $W$.

% MS lines 22–28, notion of fibration and application
Pour la notion de $W$-fibré \inserted{(\uuline{faible})}, j'ai été long à trouver enfin \uline{la} bonne notion, \deleted{\unsure{valable}} bonne pour un $W$ ne satisfaisant que $W(1,2,3)$ (et $=$, seulement $W(1,2\,\mathrm{bis},3)$, où $W(2\,\mathrm{bis})$ est l'ancien L 4). Je renvoie pour \unsure{cela} à (XII th. 3, p. 51), où le \deleted{\unsure{th.}} résultat est énoncé à \unsure{coup|coups} de condition\unsure{s} d'équivalence. En fait, pour les applications \inserted{principales} : $\mathrm{HOT}_W$, \deleted{\unsure{on}} jusqu'à présent, il suffit$^*$ d'une variante faible de $W(4\,\mathrm b)$ :

% MS bottom left, long addition across lines of main text
% slanted
\margin{C'est finalement tout \unsure{mon|un} \unsure{entier|entier,} \unsure{travail} \unsure{bien sorti|bien senti} \illegible{1} à présent. Je crois \deleted{\unsure{alors}} que \unsure{ça} \unsure{pourrait bien} \illegible{1} l'axiome $W(4\,\mathrm b')$ \unsure{strictement|strictement bien|strictement rien} \unsure{entraîner|entre-} finalement, \unsure{mais|quand} \unsure{j'aurai eu le} (à en \unsure{décider} \unsure{dans les} jours) \unsure{que} ;}

% MS last two lines, weaker condition
$W(4\,\mathrm b')$ Comme $W(4\,\mathrm b)$, mais en exigeant de plus que $X,Y$ soient dans $\underline{\mathrm{Parf}}_{\unsure{\omega|W|w}} S$.
